% !TeX program = lualatex
% Bewerbung Masterthesis – Erlind Sejdiu
% Kompilieren mit:  lualatex Anschreiben.tex
% Layout bewusst abgestimmt auf den Lebenslauf (gleiche Typografie & Farbpalette).

\documentclass[a4paper,11pt]{article}

% ─── Pakete ──────────────────────────────────────────────────────────────────
\usepackage{fontspec}
\usepackage[ngerman]{babel}
\usepackage[left=2.2cm, right=2.2cm, top=1.4cm, bottom=1.4cm]{geometry}
\usepackage{microtype}
\usepackage{enumitem}
\usepackage{xcolor}
\usepackage[hidelinks]{hyperref}

% ─── Schrift ─────────────────────────────────────────────────────────────────
\setmainfont{texgyrepagella}[
    Extension = .otf,
    UprightFont = *-regular,
    BoldFont = *-bold,
    ItalicFont = *-italic,
    BoldItalicFont = *-bolditalic
]

\definecolor{ink}{HTML}{1C1C1C}
\definecolor{muted}{HTML}{6E6E6E}
\definecolor{hairline}{HTML}{C4C4C4}
\color{ink}

\hypersetup{pdfauthor={Erlind Sejdiu},
            pdftitle={Bewerbung Masterthesis -- Erlind Sejdiu}}

\setlength{\parindent}{0pt}
\setlength{\parskip}{0.5em}
\pagestyle{empty}

% =============================================================================
\begin{document}

% ─── Kopf (identisch zum Lebenslauf) ─────────────────────────────────────────
{\LARGE Erlind Sejdiu}\\[0.3em]
{\small\color{muted}
Blücherstraße 18, Karlsruhe \enspace\textbar\enspace
0157\,3547\,4746 \enspace\textbar\enspace
\href{mailto:erlindamprogrammieren@gmail.com}{erlindamprogrammieren@gmail.com}}\\[0.4em]
\textcolor{hairline}{\rule{\linewidth}{0.7pt}}

\vspace{0.8em}

\vspace{1.1em}

% ─── Betreffzeile ────────────────────────────────────────────────────────────
{\large\textbf{Initiativbewerbung für eine Masterthesis im Bereich Software Engineering}}\\[0.15em]
{\small\color{muted}Möglicher Beginn: Sommersemester 2027 (oder nach Absprache)}

\vspace{0.5em}

% ─── Anschreiben ─────────────────────────────────────────────────────────────
Sehr geehrte Damen und Herren,% Falls Ansprechpartner bekannt: Sehr geehrte(r) Frau/Herr <Name>,

\textbf{Hier nennen wer du bist und wie du auf das Unternehmen aufmerksam geworden bist}
aktuell studiere ich Informatik im Master an der Hochschule Karlsruhe mit dem Schwerpunkt Software Engineering (aktueller Notenschnitt 1,0). Im Rahmen meines Studiums möchte ich meine Masterthesis an der Schnittstelle zwischen fundierter wissenschaftlicher Konzeption und konkreter industrieller Praxis durchführen. Die technologische Ausrichtung und der hohe Qualitätsanspruch Ihres Hauses im Bereich <konkreter Bereich / z.\,B. verteilte Systeme, moderne Cloud-Architekturen, Industrieautomation> bieten dafür das ideale Umfeld.

\textbf{Hier nimmst du die Stellenbeschreibung und die Anforderung daraus und versuchst mitzuteilen, dass das auch deine stärken sind}
Mein fachlicher Fokus liegt in der konzeptionellen Softwareentwicklung und modernen Software-Architektur. Mich begeistert die Frage, wie sich skalierbare, wartungsarme und zukunftssichere Systeme gestalten lassen. Praktische Erfahrungen bringe ich sowohl aus der Entwicklung eigener Architekturprojekte (etwa einer modularen Projektmanagement-Lösung im modernen .NET-Ökosystem mit OpenTelemetry-Instrumentierung) als auch aus meiner bisherigen industriellen Tätigkeit mit: So konnte ich bereits in meiner Bachelorarbeit (Note 1,0) sowie als Werkstudent gewachsene Softwaresysteme erfolgreich analysieren, refaktorisieren und durch automatisierte Tests stabilisieren. Fortgeschrittene Sprachfeatures, Metaprogrammierung und moderne Entwurfsmuster setze ich dabei gezielt ein, um wartbaren und performanten Code zu schaffen.



\textbf{Du musst für die Leute denken und ihnen schon Gedanken in den Kopf setzen wie man dich theoretisch einsetzen könnte (Abhängig von der Stellenbeschreibung)}
Für eine gemeinsame Masterthesis sehe ich insbesondere in drei Themenfeldern einen unmittelbaren Mehrwert für Ihr Team.
\textbf{Architekturmodernisierung \& Refactoring:} Wissenschaftlich fundierte Neukonzeption und Entkopplung gewachsener Komponenten zur Reduktion technischer Schulden.
\textbf{Prototyping innovativer Features:} Konzeption und prototypische Umsetzung zukunftsweisender Funktionen, für die im regulären Sprint-Alltag häufig Ressourcen fehlen.
\textbf{Qualität, Observability \& Robustheit:} Durchgängige Integration von Telemetrie (z.\,B. OpenTelemetry), Schärfung von Sicherheitsaspekten oder Aufbau automatisierter Teststrategien.

Über die Gelegenheit, mich und meine Motivation in einem persönlichen Gespräch vorzustellen, freue ich mich sehr.

\vspace{0.6em}

Mit freundlichen Grüßen

\vspace{1.5em}

Erlind Sejdiu

\end{document}